\chapter{Introduction}

In the current technological world, there are three different main trends:
\begin{enumerate}
    \item \ul{Computing}: computers are getting smaller, cheaper and more powerful.
    \item \ul{Sensing}: there is a huge variety of sensors that can be used to measure different physical properties. They are also getting smaller and cheaper.
    \item \ul{Transmission}: two factors are being taken into account:
    \begin{itemize}
        \item There are more and more wirelesscommunication technologies, with different characteristics. This leads to omnipresent connectivity.
        \item Energy consumption is a key factor in the design of these technologies, as many IoT devices are battery-powered. This is leading to the development of low-power communication protocols and energy-efficient hardware.
    \end{itemize}
\end{enumerate}

In this context, the Internet of Things (IoT) is emerging as a paradigm where everyday objects become smart and connected to the internet.
\begin{definicion}[IoT Thing]
    An IoT thing is a physical object that is augmented with embedded electronics, making it therefore ``smart''.
\end{definicion}

\section{Digital Twin}

The Digital Twin is a global representation of the world, where every physical object has a digital counterpart. This means that every physical object is represented in the digital world, and every change in the physical world is reflected in the digital world and vice versa. This is a very ambitious vision, and it is still far from being achieved. However, it is a useful concept to understand the potential of IoT and the challenges that it poses. Its three main features are:
\begin{itemize}
    \item \ul{Live}: the digital twin is continuously updated with data from the physical world.
    \item \ul{Predictive}: the digital twin can be used to predict the behavior of the physical world using simulations and machine learning techniques.
    \item \ul{Interactive}: depending on the predictions made by the digital twin, it can change and therefore make changes in the physical world.
\end{itemize}

In this situation, the world would become observable and programmable in real time, which would have a huge impact on many different fields, such as healthcare, transportation, manufacturing, etc. However, there are still many challenges that need to be overcome to achieve this vision.

\section{IoT Views}

There are two different views on IoT:
\begin{itemize}
    \item IoT as a network of connected \emph{things}, aka \emph{M2M (Machine to Machine)}.
    
    It is quite a low-level technical view, where the focus is on the devices and the communication between them. Internet is just seen as a communication medium, and the main goal is to connect as many things as possible.    

    It is the view that is usually taken by engineers and computer scientists. 
    \item IoT as a network of connected \emph{data}, aka \emph{Big Data}.
    
    It is a more high-level view, where there are still things but the focus is on the data that is generated by these things and how it can be used to create value. Internet is seen as a platform for data collection and analysis, and the main goal is to extract insights from the data.
    
    It is the view that is usually taken by business people and data scientists.
\end{itemize}